\chapter{Project Management and Conclusion}
\label{chap:11}

\section{Process and Schedule}
\label{sec:11.1}

The project ran across one semester as a solo build, with the supervisor setting the deliverable at a progress meeting early in it: \textit{concentrate on the justification, on why each choice was made}, with a proof of concept, a demonstration and a prototype as the evidence. Everything after that meeting was planned against that definition rather than against shipping a product.

{\small
\setlength{\LTpre}{6pt}\setlength{\LTpost}{8pt}
\begin{longtable}{>{\raggedright\arraybackslash}p{\dimexpr 0.154\textwidth-2\tabcolsep\relax}>{\raggedright\arraybackslash}p{\dimexpr 0.390\textwidth-2\tabcolsep\relax}>{\raggedright\arraybackslash}p{\dimexpr 0.456\textwidth-2\tabcolsep\relax}}
\caption{The fifteen-week plan against what was delivered}\\
\toprule
Window & Planned & Delivered \\
\midrule
\endfirsthead
\toprule
Window & Planned & Delivered \\
\midrule
\endhead
\bottomrule
\endfoot
\bottomrule
\endlastfoot
A, survey and scope & Market survey, engine shortlist, requirements, the first design decisions & Delivered. Five products surveyed, five engines shortlisted on published evidence, the fault taxonomy located in the literature \\
B, measurement & Benchmark harness, engine measured on the target device, the stack settled & Delivered. The harness is a separate application and its runs are dated and kept \\
C, build and evidence & The assessment layer, the client, the specifications, the prototype, the report & Delivered, with the engine re-decided inside this window \\
\end{longtable}
}

\textbf{The one significant re-plan happened inside window C}, and it is worth recording as a schedule event rather than only as a technical one. A second benchmark sweep, run after a frame-preparation defect in the client was fixed, changed which engine met the per-frame requirement. Acting on it meant replacing the perception layer, removing the calibration stage that the old thresholds depended on, and re-deriving every threshold from recorded video.

The cost was three report chapters that had to be written again against the system as it now is. The alternative was a report describing a design the measurements no longer supported, and 4.2 sets out why the measurement was worth acting on.

\textbf{The report itself was rebuilt once.} An earlier version, structured from the faculty template, had no natural home for either the market comparison or the justification layer, which are the two things the supervisor grades hardest. The structure here is derived from his direction rather than from the template, which is why Chapters 2, 4, 7 and 8 exist as separate chapters at all.

\section{Risk Management}
\label{sec:11.2}

{\small
\setlength{\LTpre}{6pt}\setlength{\LTpost}{8pt}
\begin{longtable}{>{\raggedright\arraybackslash}p{\dimexpr 0.327\textwidth-2\tabcolsep\relax}>{\raggedright\arraybackslash}p{\dimexpr 0.155\textwidth-2\tabcolsep\relax}>{\raggedright\arraybackslash}p{\dimexpr 0.518\textwidth-2\tabcolsep\relax}}
\caption{The risks on the register, and which ones materialised}\\
\toprule
Risk & Materialised & Outcome \\
\midrule
\endfirsthead
\toprule
Risk & Materialised & Outcome \\
\midrule
\endhead
\bottomrule
\endfoot
\bottomrule
\endlastfoot
The handset cannot sustain the frame rate & No & 29.57 fps at the camera ceiling, with 2.4 times the latency bound in hand \\
No labelled training data exists for this project & Yes, as expected & Answered by design: geometric constraints rather than a learned classifier, 4.5 \\
Published sources give criteria but not thresholds & Yes & Every threshold measured and labelled as measured; the provenance field makes the gap visible rather than hidden \\
A fault is not observable in the declared camera view & Yes, twice & Each removed rather than approximated, with the reason in the specification file \\
Thermal throttling over a long session & Yes, and bounded & The device throttles inside the first three minutes and then holds: 25.00 fps at the worst of twenty windows, against a floor of 22 (9.4.1) \\
A defect invisible to the test suite & Yes, three times & Found on the device and in recorded video, which is why the replay harness exists \\
\end{longtable}
}

Two of these are worth more than a table row.

\textbf{The absence of labelled data was the binding constraint on the whole design}, and it was identified at the start rather than discovered. It is what removes three of the four assessment approaches in 4.5, and the declarative format in 4.6 exists because of it. A risk that shapes the architecture is better handled early than mitigated late.

\textbf{The defects that unit tests could not see} were the most expensive class of problem in the project, and all three shared a shape: the evaluator behaved exactly as specified, and the specification described something the movement does not do. A threshold set against one engine's noise, a repetition bound set from anatomy rather than from the measurement, and a trigger control counted in the wrong unit. None is a coding error and none would fail a test. The offline replay harness was built in response.

\section{What Was Delivered}
\label{sec:11.3}

Against the objectives in 1.3.

\textbf{O1, on-device perception. Delivered.} MoveNet Lightning on the GPU delegate, 29.57 fps and 27 ms at the 95th percentile on the target handset, against requirements of 22 fps and 65 ms. No frame leaves the device and none is written to storage.

\textbf{O2, movement assessment. Delivered.} Three exercises across two movement categories, six faults, one generic evaluator with no branch on any exercise name. Repetition counting exact on all seven counting sets recorded.

\textbf{O3, immediate spoken feedback. Delivered.} A pre-authored cue bank with no language model in the live path. Confirmation to audio dispatch measured at a median of 16 ms on the device.

\textbf{O4, the application, a server service and a user profile. Partially delivered.} The application exists and records sessions locally. The server is specified, its contract is fixed by FR-16 to FR-18, and it is not built in this phase. This is the one objective not met in full, and 11.4 states it as a divergence rather than leaving it to be inferred from a missing chapter.

\textbf{O5, traceable fault definitions and a labelling protocol. Delivered.} Five of the six faults carry a cited criterion and all six carry a provenance record, enforced by the format rather than by convention. The blind labelling protocol is written, and a certified trainer has reviewed the recordings and confirmed that each performs the fault its set is named for. The blind pass itself, which is the stronger claim, has not been run.

\textbf{O6, a labelled dataset and a testing procedure. Delivered.} Ten scripted sets, 9,893 frames, with the condition known by construction and the repetition count confirmed independently of the system. The measurement procedure is written in enough detail to be repeated, and the replay harness that scores it runs the shipped model.

\section{Divergences from the Registration}
\label{sec:11.4}

Three, each named rather than absorbed.

\textbf{The server is specified and not built.} The registration describes a backend service with user accounts. The contract exists as requirements and the client is written to produce exactly what it would consume, so the boundary is real and the work is bounded. What does not exist is the service.

\textbf{Movement assessment is geometric rather than learned.} The registration anticipated a model trained to classify technique faults. 4.5 sets out why: such a model needs hundreds to thousands of expert-labelled examples per fault per exercise, this project has none, and inventing labels to train against would have made every figure in Chapter 9 meaningless. The declarative layer is designed so that the evaluator behind a fault code can later be learned without changing the specification format or the cue path.

\textbf{The dataset is methodology plus a scripted corpus rather than a large labelled one.} Ten sets filmed to a written protocol, with the independent expert pass specified and agreed. The registration's ambition here was larger, and what makes this version defensible is that the procedure is written down precisely enough for the corpus to be extended by someone else.

The registration itself is unchanged, on the supervisor's explicit direction that it stay ambitious across both courses.

\section{Phase 2 Scope}
\label{sec:11.5}

1.4.2 set the boundary. This phase reads the movement. The next one explains it. What follows is the direction of that phase. It is not a plan, and it fixes no order and no dates.

The assessment layer is what the rest of the product is built on, and it is what this phase delivered. Everything that comes next is added next to it, not inside it. A coaching layer that explains a fault in words between sets. A server that keeps sessions across devices. Further features that improve the experience of using the app. 4.4 keeps language generation out of the live feedback path, which is why these can be added without changing the architecture.

Which of them comes first is decided at the start of that phase, not here.

\section{Conclusion}
\label{sec:11.6}

This phase set out to establish whether a phone can usefully correct a movement while the movement is still happening. The measurements support yes for this phase, within the limits 9.5 sets out.

The perception layer holds 29.57 frames per second at 27 milliseconds for the 95th percentile of its per-frame path, on a handset the project already had, with 2.4 times the required latency margin. Across ten recorded sets and 9,893 frames it did not lose the body once.

The assessment layer counted every repetition correctly across seven sets and detected every scripted fault. Precision was 0.990 across the five faults judged frame by frame, and squat depth, which is judged once per repetition, was raised on 10 of 10 shallow repetitions and none of the 10 clean ones. No fault fell below 0.957 against a requirement of 0.90. The two faults most at risk of being confused with one another, a hip dropping and a hip rising, were separated completely in both exercises that carry them.

Three properties of how that was done matter more than the figures.

\textbf{Every cue traces to a published criterion}, and where it does not, the specification says so in its own fields rather than in prose a reader has to find. Five of six faults carry a source. No threshold does, and that is reported as a result.

\textbf{Every number in this report names the run that produced it.} The two requirements that were not measured are stated as not measured, in the chapter that would otherwise have claimed them.

\textbf{An exercise is a file, not a release.} Three exercises across two movement categories run through one evaluator containing no code specific to any of them, which is checked by the suite rather than asserted here.

What the phase does not establish is equally clear. One subject, one session, one handset. Thresholds measured from the performer they were later evaluated on. An expert labelling pass agreed and not yet performed. Those are the boundary of the claim, and they are the work the second phase starts from.
